\section{Síntesis y ampliación}

\begin{enumerate}
    \item Los problemas abiertos del RL profundo abarcan tres niveles: la definición del problema, los métodos y el despliegue
    \item El diseño de recompensas (para evitar el reward hacking), la capacidad de generalización y la seguridad son las líneas más importantes
    \item La buena investigación parte de entender el problema, no del método
    \item Los experimentos de RL requieren varias semillas aleatorias, líneas base sólidas y experimentos de ablación
    \item La depuración de RL requiere un método sistemático; la visualización y la verificación en entornos simples son esenciales
    \item La integración entre áreas (RL + LLM, RL + robotics) es la frontera más activa en la actualidad
\end{enumerate}

\begin{knowledgebox}{Lecturas adicionales}
\begin{itemize}
    \item Henderson et al., ``Deep Reinforcement Learning that Matters,'' AAAI 2018
    \item Amodei et al., ``Concrete Problems in AI Safety,'' arXiv 2016
    \item Engstrom et al., ``Implementation Matters in Deep Policy Gradients,'' ICLR 2020
    \item Andrychowicz et al., ``What Matters in On-Policy Reinforcement Learning? A Large-Scale Empirical Study,'' ICLR 2021
\end{itemize}
\end{knowledgebox}

\end{document}
